% Generated from locale/tr/content/first-order-logic/syntax-and-semantics/terms-formulas.tex; do not hand-edit.


\olfileid[tr]{fol}{syn}{frm}

\olsection{Terimler ve \printtoken{P}{formula}}

Birinci dereceden bir dil~$\Lang L$ verildiğinde, $\Lang L$'nin temel söz
varlığından kurulan ifadeleri tanımlayabiliriz. Bunlar arasında özellikle
\emph{terimler} ve \emph{!!{formula}s} bulunur.

\begin{defn}[Terimler]
\ollabel{defn:terms}
$\Lang L$'nin \emph{terimleri} kümesi~$\Trm[L]$ tümevarımsal olarak
şöyle tanımlanır:
\begin{enumerate}
\item Her !!{variable} bir terimdir.
\item $\Lang L$'nin her !!{constant}{}i bir terimdir.
\item $f$, $n$-yerli bir !!{function} ve $t_1$, \dots, $t_n$
  terimlerse $\Atom{f}{t_1, \ldots, t_n}$ bir terimdir.
\tagitem{limitClause}{
  Başka hiçbir şey terim değildir.}{}
\end{enumerate}
İçinde !!{variable}s bulunmayan terime \emph{kapalı terim} denir.
\end{defn}

\begin{explain}
!!{constant}s, dili ve terimleri belirtirken ayrı bir simge kategorisi
olarak yer alır; oysa bunları sıfır yerli !!{function}s arasında da
sayabilirdik. Böylece terim tanımının ikinci maddesine gerek kalmazdı.
Yalnızca $n = 0$ olduğunda $\Atom{f}{t_1, \ldots, t_n}$ ifadesini tek
başına $f$ olarak anlamamız gerekirdi.
\end{explain}

\begin{defn}[Formüller]
\ollabel{defn:formulas}
$\Lang L$ dilinin \emph{!!{formula}s} kümesi~$\Frm[L]$ tümevarımsal
olarak şöyle tanımlanır:
\begin{enumerate}
\tagitem{prvFalse}{$\lfalse$ atomik bir !!{formula}{}dür.}{}

\tagitem{prvTrue}{$\ltrue$ atomik bir !!{formula}{}dür.}{}

\item $R$, $\Lang L$'nin $n$-yerli bir !!{predicate}{}i ve $t_1$, \dots,
  $t_n$, $\Lang L$'nin terimleriyse $\Atom{R}{t_1,\ldots, t_n}$ atomik bir
  !!{formula}{}dür.

\item $t_1$ ve $t_2$, $\Lang L$'nin terimleriyse $\Atom{\eq}{t_1, t_2}$
  atomik bir !!{formula}{}dür.

\tagitem{prvNot}{$!A$ bir !!a{formula} ise $\lnot !A$ da bir
  !!a{formula}{}dür.}{}

\tagitem{prvAnd}{$!A$ ve $!B$ !!{formula}s ise $(!A \land
  !B)$ bir !!a{formula}{}dür.}{}

\tagitem{prvOr}{$!A$ ve $!B$ !!{formula}s ise $(!A \lor !B)$
  bir !!a{formula}{}dür.}{}

\tagitem{prvIf}{$!A$ ve $!B$ !!{formula}s ise $(!A \lif !B)$
  bir !!a{formula}{}dür.}{}

\tagitem{prvIff}{$!A$ ve $!B$ !!{formula}s ise $(!A \liff !B)$
  bir !!a{formula}{}dür.}{}

\tagitem{prvAll}{$!A$ bir !!a{formula} ve $x$ bir !!a{variable} ise
  $\lforall[x][!A]$ bir !!a{formula}{}dür.}{}

\tagitem{prvEx}{$!A$ bir !!a{formula} ve $x$ bir !!a{variable} ise
  $\lexists[x][!A]$ bir !!a{formula}{}dür.}{}

\tagitem{limitClause}{Başka hiçbir şey bir !!a{formula} değildir.}{}
\end{enumerate}
\end{defn}

\begin{explain}
Terimler kümesinin ve !!{formula}s kümesinin tanımları
\emph{tümevarımsal tanımlar}dır. Özünde !!{formula}s kümesini sonsuz
sayıda aşamada kurarız. İlk aşamada bütün atomik formülleri formül ilan
ederiz; bu, tanımın ilk birkaç durumuna, yani
\iftag{prvTrue}{$\ltrue$, }{}%
\iftag{prvFalse}{$\lfalse$, }{}%
$\Atom{R}{t_1,\dots,t_n}$ ve $\Atom{\eq}{t_1,t_2}$ durumlarına karşılık
gelir. Dolayısıyla ``atomik !!{formula}'', bu biçimlerden herhangi birindeki
!!{formula} demektir.

Tanımın öteki durumları, daha önce kurulmuş !!{formula}s{}den yeni
!!{formula}s kurma kuralları verir. İkinci aşamada bu kurallarla atomik
!!{formula}s{}den formüller kurabiliriz. Üçüncü aşamada atomik
!!{formula}s{}den ve ikinci aşamada elde edilenlerden yeni formüller kurarız;
bu böyle sürer. Bir aşamada sonunda kurulabilen her şey !!{formula}{}dür
ve başka hiçbir şey formül değildir.
\end{explain}

Uzlaşım gereği $\eq$ simgesini argümanlarının arasına yazar ve
parantezleri kaldırırız: $\eq[t_1][t_2]$,
$\Atom{\eq}{t_1,t_2}$ için bir kısaltmadır. Ayrıca
$\lnot \Atom{\eq}{t_1,t_2}$, $\eq/[t_1][t_2]$ biçiminde kısaltılır.
$!B$ ve $!C$'den iki yerli bir bağlaç~$\ast$ kullanılarak kurulan
$(!B \ast !C)$ formülünü yazarken çoğu kez en dıştaki parantez çiftini
atar ve yalnızca~$!B \ast !C$ yazarız.

\begin{intro}
Bazı mantık metinleri, !!{variable}~$x$'in $!A$'da geçmesini,
\iftag{prvEx}{$\lexists[x][!A]$ }{}%
\iftag{notprvEx,notprvAll}{}{ve }%
\iftag{prvAll}{$\lforall[x][!A]$ }{}%
ifadesinin
\iftag{notprvEx,notprvAll}{bir !!a{formula}}{!!{formula}s}
sayılması için şart koşar.
Bunu şart koşmazsanız kötü bir şey olmaz; üstelik işler kolaylaşır.
\end{intro}

\begin{tagblock}{defNot,defOr,defAnd,defIf,defIff,defTrue,defFalse,defEx,defAll}
\begin{defn}
Tanımlı işleçlerle kurulan formüller şöyle anlaşılmalıdır:

\begin{tagenumerate}{defTrue,defFalse,defNot,defOr,defAnd,defIf,defIff,defEx,defAll}

\tagitem{defTrue}{$\ltrue$,
  \iftag{prvFalse}{$\lnot\lfalse$}{sabitlenmiş atomik bir
    !!{formula}~$!A$ için $(!A \lor \lnot !A)$} ifadesini kısaltır.}{}

\tagitem{defFalse}{$\lfalse$,
  \iftag{prvTrue}{$\lnot\ltrue$}{sabitlenmiş atomik bir
    !!{formula}~$!A$ için $(!A \land \lnot !A)$} ifadesini kısaltır.}{}

\tagitem{defNot}{$\lnot !A$, $!A \lif \lfalse$ ifadesini kısaltır.}{}

\tagitem{defOr}{$!A \lor !B$,
  \iftag{prvAnd}{$\lnot(\lnot !A \land \lnot !B)$}{$\lnot !A \lif
    !B$} ifadesini kısaltır.}{}

\tagitem{defAnd}{$!A \land !B$,
  \iftag{prvOr}{$\lnot(\lnot !A \lor \lnot !B)$}{$\lnot (!A \lif
    \lnot !B)$} ifadesini kısaltır.}{}

\tagitem{defIf}{$!A \lif !B$,
  \iftag{prvOr}{$\lnot !A \lor !B$}{$\lnot (!A \land \lnot !B)$}
  ifadesini kısaltır.}{}

\tagitem{defIff}{$!A \liff !B$, $(!A \lif !B) \land (!B
  \lif !A)$ ifadesini kısaltır.}{}

\tagitem{defAll}{$\lforall[x][!A]$, $\lnot\lexists[x][\lnot !A]$
  ifadesini kısaltır.}{}

\tagitem{defEx}{$\lexists[x][!A]$, $\lnot\lforall[x][\lnot !A]$
  ifadesini kısaltır.}{}
\end{tagenumerate}
\end{defn}
\end{tagblock}

Belirli bir uygulamanın diliyle çalışıyorsak, iki yerli !!{predicate}s
ve !!{function}s söz konusu olduğunda bunları çoğu kez ilgili terimlerin arasına yazarız; örneğin
aritmetiğin dilinde $t_1 < t_2$ ile $(t_1 + t_2)$, küme teorisinin
dilinde ise $t_1 \in t_2$. Aritmetiğin dilindeki ardıl fonksiyonu,
uzlaşım gereği argümanından \emph{sonra} bile yazılır:~$t'$. Ne var ki
resmî olarak bunlar sırasıyla $\Atom{\Obj A^2_0}{t_1, t_2}$,
$\Obj f^2_0(t_1, t_2)$, $\Atom{\Obj A^2_0}{t_1, t_2}$ ve $f^1_0(t)$
için yalnızca uzlaşımsal kısaltmalardır.

\begin{defn}[Sözdizimsel özdeşlik]
$\ident$ simgesi, simge dizileri arasındaki sözdizimsel özdeşliği ifade
eder; başka bir deyişle $!A \ident !B$ ancak ve ancak $!A$ ile $!B$ aynı
uzunlukta simge dizileriyse ve her konumda aynı simgeyi içeriyorsa.
\end{defn}

$\ident$ simgesinin iki yanında art arda eklemeyle elde edilmiş simge
dizileri bulunabilir. Örneğin $!A \ident (!B \lor !C)$, $!A$ simge
dizisinin sırasıyla bir açma parantezi, $!B$ dizisi, $\lor$ simgesi,
$!C$ dizisi ve bir kapama parantezi art arda eklenerek elde edilen
diziyle aynı olduğu anlamına gelir. Durum buysa $!A$'nın ilk simgesinin
bir açma parantezi olduğunu, $!A$'nın $!B$'yi (ikinci simgeden
başlayarak) bir alt dizi olarak içerdiğini, bu alt diziyi~$\lor$'un
izlediğini vb. biliriz.

Terimler ve !!{formula}s, temel öğelerden tümevarımsal tanımlarla
kurulduğu için bunlara ilişkin özellikleri ispatlamak üzere aşağıdaki
tümevarım ilkelerini kullanabiliriz.

\begin{lem}[\emph{Terimler üzerinde tümevarım ilkesi}]
    \ollabel{lem:trmind}
    $\Lang L$ birinci dereceden bir dil olsun.
    Bir $P$ özelliği
    %
    \begin{enumerate}
        \item her !!{variable}~$v$ için geçiyorsa,
        %
        \item $\Lang L$'nin her !!{constant}~$a$'sı için geçiyorsa ve
        %
        \item $f$, $\Lang L$'nin $n$-yerli bir !!{function}{}u olmak
        üzere bu özellik $t_1$,~\dots, $t_n$ için geçtiğinde
        $f(t_1,\dotsc,t_n)$ için de geçiyorsa
    \end{enumerate}
    ($t_1$,~\dots, $t_n$'nin $\Lang{L}$ terimleri olduğu varsayımıyla),
    $P$, $\Trm[L]$'deki her terim için geçer.
\end{lem}

\begin{prob}
    \olref[fol][syn][frm]{lem:trmind}'yı ispatlayın.
\end{prob}

\begin{lem}[\emph{!!{formula}s üzerinde tümevarım ilkesi}]
    \ollabel{thm:frmind}
    $\Lang L$ birinci dereceden bir dil olsun.
    Bir $P$ özelliği bütün atomik !!{formula}s için geçiyorsa ve
    %
    \begin{enumerate}
        \tagitem{prvNot}{özellik $!A$ için geçtiğinde $\lnot !A$ için de
            geçiyorsa; }{}
        \tagitem{prvAnd}{özellik $!A$ ve~$!B$ için geçtiğinde
            $(!A \land !B)$ için de geçiyorsa; }{}
        \tagitem{prvOr}{özellik $!A$ ve~$!B$ için geçtiğinde
            $(!A \lor !B)$ için de geçiyorsa; }{}
        \tagitem{prvIf}{özellik $!A$ ve~$!B$ için geçtiğinde
            $(!A \lif !B)$ için de geçiyorsa; }{}
        \tagitem{prvIff}{özellik $!A$ ve~$!B$ için geçtiğinde
            $(!A \liff !B)$ için de geçiyorsa; }{}
        \tagitem{prvEx}{özellik $!A$ için geçtiğinde
            $\lexists[x][!A]$ için de geçiyorsa; }{}
        \tagitem{prvAll}{özellik $!A$ için geçtiğinde
            $\lforall[x][!A]$ için de geçiyorsa; }{}
  \end{enumerate}
  ($!A$ ile $!B$'nin $\Lang{L}$'nin !!{formula}s olduğu varsayımıyla),
  $P$, $\Frm[L]$'deki bütün formüller için geçer.
\end{lem}

